\section{The model}\label{sec:02_model}

\subsection{The economic setting}\label{sec:02_model-setting}

The model fits a listed target in a \emph{decision interval} \paper{Sec.~2.1}. In this interval:
\begin{itemize}
\item the public knows about the sale opportunity (for example, a disclosed approach, an announced
strategic review or an open contest);
\item the stock still trades;
\item one buyer is prepared, and a further buyer can still come in.
\end{itemize}
Preparation means due diligence, financing, approvals and the use of evaluation capacity. The cost is sunk.
A buyer that does not prepare cannot bid its expected value; it stays out of the sale. So the cost is an
\emph{entry cost}: payment reveals the buyer's value and permits a bid. It is not the price of an optional
report.

Incumbent strength is a \emph{distribution}, not a bid. The public record can identify a strategic buyer with clear synergies.
It can instead identify a financial buyer with high leverage. It
need not show what either buyer would pay. The model varies the public distribution of the incumbent's
value and keeps the realized value private until the bid.

\subsection{Players, values and information}\label{sec:02_model-players}

Table~\ref{tab:02_model-players} lists the players.

\begin{table}[!ht]
\centering\small
\begin{tabularx}{\linewidth}{@{}>{\bfseries}l X X@{}}
\toprule
\textnormal{Player} & Primitive & What it knows when it acts\\
\midrule
Incumbent & Value \(R\sim F\), continuous on \([0,\bar r]\); benchmark \(F\) uniform on \([0,r]\). Already
prepared, no further cost. & Its own \(R\) (at the bid stage).\\[3pt]
Challenger & Class \(j\in\{H,L\}\) with \(\Pr(j=H)=\pi\); value \(\theta\sim G_j\) given the class;
benchmark \(\theta=h\) in class \(H\) and \(\theta=\ell\) in class \(L\), \(\pi=1/2\). Cost \(C=0\) with
probability \(\rho\), else \(c_H\). & The price \(P\); then its own cost \(C\). Not \(j\), \(\theta\),
\(X\) or \(R\). It learns \(\theta\) only if it prepares.\\[3pt]
Investor & Fills the slot with probability \(\lambda\). Order \(q\in\{-1,0,1\}\); cost \(k\) for a nonzero
order. No position in the target, no control rights. & The class \(j\) (not \(\theta\) in general; in the
benchmark the class fixes \(\theta\)). Only the distribution of \(R\).\\[3pt]
Liquidity trader & Fills the slot with probability \(1-\lambda\). Sells, abstains or buys one unit with
probabilities \(\beta\), \(1-2\beta\), \(\beta\); \(0<\beta\le1/3\). & Nothing about the sale (it trades
for reasons apart from the sale).\\[3pt]
Market makers & Competitive and risk neutral. Set \(P(x)=\E[V_T\mid X=x]\) and absorb the order. & The
order \(X\in\{-1,0,1\}\) only, not who placed it.\\[3pt]
Seller & Commits to a cash second-price auction with reserve \(p>0\). & Moves first.\\
\bottomrule
\end{tabularx}
\caption{Players, primitives and information \paper{Secs.~2.2--2.3}.}
\label{tab:02_model-players}
\end{table}

The classes satisfy the \emph{class ranking}:
\[
G_H(y)\le G_L(y)\ \text{for every }y,\qquad G_H(p^-)=G_L(p^-)=0. \tag*{\paper{(1)}}
\]
The first part is first-order stochastic dominance: class \(H\) has weakly higher values. The second part puts
both class supports at or above the reserve, so every value meets it. The difference
\(D(y)=G_L(y)-G_H(y)\ge0\) recurs below. In the benchmark, \(G_H\) is a point mass at \(h\) and \(G_L\)
is a point mass at \(\ell\), so the class and the value coincide and \(D(y)=\1\{\ell\le y<h\}\).

The benchmark ordering is
\[
0<p<\ell<r<h. \tag*{\paper{(2)}}
\]
Read it as follows:
\begin{itemize}
\item \(p<\ell\): the reserve excludes neither class;
\item \(\ell<r\): the incumbent can beat a low-value challenger, with probability \((r-\ell)/r\);
\item \(r<h\): the incumbent never beats a high-value challenger.
\end{itemize}

The pair \((j,\theta)\), the incumbent's value \(R\), the cost \(C\) and the draws of the order slot are
mutually independent. So the cost carries no information about the class or the value, and the order
carries no information about \(\theta\) beyond the class. The standalone value of the target is zero, and
all values and costs are per target share. A common rescale of every value, the reserve and every cost
rescales prices and payoffs; it leaves orders, beliefs, entry and ownership unchanged. So the guide can
measure money in units of \(\ell\).

\begin{panelnote}
\textbf{``Why does the challenger not know its own class?''} Value depends on facts about the target:
its customers, its technology, the demand for its product. A specialist investor can know these facts
better before diligence. A buyer knows its own integration plans, which are different facts about the
same match. Diligence reveals \(\theta\). The benchmark gives the investor exact knowledge of the class
and the buyer none, so that the mechanism is easy to see. Section~6.3 of the paper
(\Cref{sec:10_robustness-signals} here) gives both sides private signals about different facts. In an
example the buyer's signal is \emph{more} accurate than the investor's, and the rise in entry needs both
signals. Section~5 of the paper (\Cref{sec:09_research} here) makes the investor's information the
result of costly research.
\end{panelnote}

\subsection{Timing}\label{sec:02_model-timing}

\begin{figure}[htbp]
\centering
\begin{tikzpicture}[font=\small]
\draw[-{Stealth[length=2.5mm]}, thick] (0,0) -- (16.2,0);
\foreach \x/\n in {1.3/1, 4.0/2, 6.7/3, 9.4/4, 12.1/5, 14.8/6} {
  \fill[intuit] (\x,0) circle (2pt);
  \node[above] at (\x,0.08) {\n};
}
\node[below, text width=2.45cm, align=center] at (1.3,-0.12) {Seller announces a cash second-price auction with reserve \(p\)};
\node[below, text width=2.45cm, align=center] at (4.0,-0.12) {Investor learns the class \(j\)};
\node[below, text width=2.45cm, align=center] at (6.7,-0.12) {Slot draw: the investor (prob.\ \(\lambda\)) or a liquidity trader places \(X\in\{-1,0,1\}\)};
\node[below, text width=2.45cm, align=center] at (9.4,-0.12) {Market makers see \(X\) and set \(P=\E[V_T\mid X]\)};
\node[below, text width=2.45cm, align=center] at (12.1,-0.12) {Challenger sees \(P\), then learns its cost \(C\); decides whether to prepare};
\node[below, text width=2.45cm, align=center] at (14.8,-0.12) {Prepared bidders learn values and bid truthfully; payoffs realized};
\end{tikzpicture}
\caption{The sequence of moves \paper{Sec.~2.3}.}
\label{fig:02_model-timing}
\end{figure}

Three features matter most:
\begin{enumerate}
\item The price forms \emph{before} the challenger decides. So the price can change the decision.
\item The market makers know this. They price the entry that each price will cause. This is the
\emph{feedback} from the financial market to a real decision.
\item No one sees who filled the slot. A liquidity order can match the investor's direction, so no order
is conclusive about the class.
\end{enumerate}

\subsection{Payoffs and the trading stage}\label{sec:02_model-payoffs}

\paragraph{The investor's profit.} If the investor fills the slot with order \(q\), its order is the whole
order that market makers see. It trades at the price \(P(q)\) and earns
\[
q\,[V_T-P(q)]-k|q|. \tag*{\paper{(3)}}
\]
Here \(V_T\) is the terminal payoff of one target share: what the auction pays per share, or zero if there
is no sale. The cost \(k\) is a separate friction (an execution cost, a cost to carry the position, the
cost to borrow shares for a short sale). It is \emph{not} price impact. Price impact comes from the
competitive rule \(P(x)=\E[V_T\mid X=x]\).

\paragraph{The slot parameters.} Define
\[
b=(1-\lambda)\beta,\qquad b_0=(1-\lambda)(1-2\beta),\qquad u=\lambda/b.
\]
Here \(b\) is the liquidity mass in each direction, and \(b_0\) is the liquidity mass at zero.
The ratio \(u\) compares informed presence with liquidity in one direction.
\Cref{sec:04_slot_beliefs-bounds,sec:04_slot_beliefs-zero} derive the bounds
\[
m(\pi)=\frac{\pi}{\pi+(1-\pi)(1+u)},\qquad
M(\pi)=\frac{\pi(1+u)}{\pi(1+u)+1-\pi}.\tag*{\paper{(4)}}
\]
A larger \(u\) moves the two bounds apart: each order says more about the class.
At the prior \(1/2\) the two bounds are symmetric about \(1/2\).
\Cref{sec:benchmark-slot} gives the benchmark slot.

\paragraph{Expected payoff given the class and the order.}
Let \(e(q)\) be preparation after the price \(P(q)\).
Given the class, preparation is independent of both acquisition values.
Thus the expected terminal payoff is \(t_0+e(q)(t_j-t_0)\), and the investor maximizes
\[
U_j(q)=q\,[t_0+e(q)(t_j-t_0)-P(q)]-k|q|.
\]
The ex ante payoff is \(\lambda U_j(q)\); this positive factor leaves the best order unchanged.
\Cref{sec:04_slot_beliefs-price} derives the conditional payoff from the two preparation outcomes.

\paragraph{The competitive price.} Average the terminal payoff over the class at posterior \(\mu(x)\):
\[
P(x)=t_0+e(x)\bigl[t_L-t_0+\DT\mu(x)\bigr].\tag*{\paper{(5)}}
\]
\Cref{sec:04_slot_beliefs-price} gives each algebra step.
\Cref{sec:05_investor-residual} then derives the investor's residual advantage after this price.

\begin{intuition}[Why the model uses these devices]
\begin{itemize}
\item \textbf{Alternative arrival.} Either trader fills one slot, as in Dow and Gorton (1997).
Every class can reach every order through liquidity trade.
Thus no order reveals the class with certainty.
\item \textbf{The unit menu.} Every order has positive probability under every profile.
Bayes' rule therefore applies without off-path beliefs.
\item \textbf{The restriction \(\beta\le1/3\).}
It extends the posterior bound to zero and gives the universal ceiling.
Propositions~2 and~3 need only \(\beta<1/2\).
\item \textbf{The cost \(k>0\).}
A small residual advantage need not cover this cost.
This permits no informed trade against a weak incumbent.
\item \textbf{Zero preparation cost with probability \(\rho\).}
It supplies entry after every price.
With \(\DT>0\), this floor permits price inversion.
It also supplies a lower bound for correctly signed nonzero orders.
Only the high-cost challenger responds to the price.
\end{itemize}
\end{intuition}

\subsection{Equilibrium}\label{sec:02_model-equilibrium}

An equilibrium specifies orders, beliefs, prices, preparation, and bids \paper{Sec.~2.4}.
Each class can mix over the three orders, including wrong-signed orders.
The market makers see the order; the challenger sees the price.

Formally \paper{Def.~A.1}, an equilibrium is a tuple \(\mathcal E=(\sigma_H,\sigma_L,\mu,P,\bar\mu,\zeta)\)
together with the bids. Three objects recur:
\begin{itemize}
\item the liquidity order law \(\eta(-1)=\eta(1)=\beta\), \(\eta(0)=1-2\beta\);
\item the probability of the order \(x\) in class \(j\), \(f_j(x)=(1-\lambda)\eta(x)+\lambda\sigma_j(x)\);
\item entry after the order \(x\), \(e(x)=\rho+(1-\rho)\zeta(P(x))\), where \(\zeta(\bar P)\) is the
probability that a high-cost challenger prepares after the price \(\bar P\).
\end{itemize}
The tuple satisfies six conditions. Read them one by one.
\begin{enumerate}[label=(E\arabic*),leftmargin=3em]
\item \textbf{Bayes consistency.} \(\mu(x)=\pi f_H(x)/[\pi f_H(x)+(1-\pi)f_L(x)]\) for every order. The
market makers see the order, not who placed it. Since \(f_j(x)\ge(1-\lambda)\eta(x)>0\), the formula is
defined at every order.
\item \textbf{Challenger beliefs.} The challenger sees the price, not the order. If several orders share
one price, it pools them: \(\bar\mu(\bar P)\) is the posterior of class \(H\) given the event
\(\{x:P(x)=\bar P\}\). Its odds are \(\pi\sum f_H(x)/[(1-\pi)\sum f_L(x)]\), with both sums over that
event.
\item \textbf{Entry optimality.} \(\zeta(\bar P)=1\) if \(B(\bar\mu(\bar P))>c_H\), and \(\zeta(\bar P)=0\)
if \(B(\bar\mu(\bar P))<c_H\). At equality, the tie rule applies. The zero-cost challenger prepares after
every price; its gross profit is an expected auction rent, which is never negative.
\item \textbf{Competitive prices.} \(P(x)=\E[V_T\mid X=x]\). Market makers earn zero in expectation after
each order. With (E1) and the independence of the cost, this condition gives \paper{(5)}.
\item \textbf{Investor optimality.} \(\sigma_j(q)>0\) only if \(q\) maximizes \(U_j(q)\) over
\(\{-1,0,1\}\), with \(P(\cdot)\) and \(e(\cdot)\) held at their values in \(\mathcal E\).
\item \textbf{Truthful bids.} Each prepared bidder whose value meets the reserve bids its value. This is
weakly dominant (\Cref{sec:03_auction-rules}).
\end{enumerate}
The appendix of the paper uses two words. A \emph{candidate} is a tuple that satisfies (E1)--(E4) and (E6)
for some profile. An \emph{equilibrium} is a candidate that also satisfies (E5). The uniqueness proofs derive bounds for every candidate before they verify the proposed tuple.

\begin{caution}[two different comparisons]
\begin{itemize}
\item A \textbf{unilateral deviation} by the investor changes only its own order. The price function
\(P\) and the preparation rule \(\zeta\) \emph{stay fixed}. The deviation payoff is \(U_j(q)\) read at
the candidate's \(P(q)\) and \(e(q)\). Use this to test whether a candidate is an equilibrium.
\item A \textbf{comparison across economies} (for example \(r_0\) against \(r_1\)) \emph{solves again}
for the prices and the preparation rule in each economy. Use this to compare outcomes.
\end{itemize}
Do not mix the two on the panel. A deviation never moves the price schedule. Because every order is on
path, a deviation always reaches a price that the candidate already defines.
\end{caution}

\paragraph{The default tie conventions.} Two indifferences can arise.
\begin{itemize}
\item \textbf{Rule IM} (investor): an indifferent investor may mix over its best orders. The alternative
\textbf{rule I0} selects abstention when abstention ties with a nonzero order.
\item \textbf{Rule C1} (challenger): an indifferent high-cost challenger prepares. The alternative
\textbf{rule CM} lets it prepare with any probability.
\end{itemize}
IM and C1 are the defaults \paper{Sec.~2.4}. Proposition~2 does not depend on them: every comparison in
its proof is strict, so the unique outcome is the same under every tie rule. The tie rules matter on the
coexistence path (\Cref{sec:08_coexistence}). The threshold branch in the paper's Figure~2 needs CM, except at its left endpoint \(r_A\), which also satisfies C1.

\subsection{Notation that recurs}\label{sec:02_model-notation}

Table~\ref{tab:02_model-notation} collects the symbols that recur in the guide.

\begin{table}[!ht]
\centering\small
\begin{tabularx}{\linewidth}{@{}l X@{}}
\toprule
Symbol & Meaning\\
\midrule
\(F\), \(G_j\) & incumbent value law; value law of class \(j\in\{H,L\}\)\\
\(D(y)=G_L(y)-G_H(y)\) & class difference at the value level \(y\); \(D\ge0\) under the class ranking\\
\(t(\theta)\), \(g(\theta)\) & expected target proceeds and challenger gross profit at a known value \(\theta\)\\
\(t_0\) & expected target proceeds if the challenger does not prepare\\
\(t_H\), \(t_L\) & expected target proceeds if a class-\(H\) / class-\(L\) challenger prepares\\
\(g_H\), \(g_L\) & expected gross profit of a prepared class-\(H\) / class-\(L\) challenger\\
\(\DT=t_H-t_L\) & target-payoff spread\\
\(w_j=t_j-t_0\), \(w(\mu)\) & entry gain in class \(j\); its average \(t_L-t_0+\DT\mu\) at belief \(\mu\)\\
\(\xi(R)\) & expected class difference in proceeds at incumbent value \(R\); \(\DT=\E_F[\xi(R)]\)\\
\(B(\mu)=g_L+\mu(g_H-g_L)\) & gross profit at belief \(\mu\); \(B_F\), \(B_r\) when the incumbent matters\\
\(\tau=(c_H-g_L)/(g_H-g_L)\) & preparation threshold, defined when \(g_H>g_L\): \(B(\mu)\ge c_H\) if and only if \(\mu\ge\tau\)\\
\(\lambda\), \(\beta\), \(\eta\) & presence probability; liquidity probability; liquidity order law\\
\(b\), \(b_0\), \(u\) & \((1-\lambda)\beta\); \((1-\lambda)(1-2\beta)\); \(\lambda/b\)\\
\(\sigma_j\), \(f_j(x)\) & investor order law in class \(j\); probability of order \(x\) in class \(j\)\\
\(\mu(x)\), \(\bar\mu(\bar P)\) & market makers' posterior after order \(x\); challenger's belief at price \(\bar P\)\\
\(m(\pi)\), \(M(\pi)\) & bottom and top posterior \paper{(4)}\\
\(\kappa(\pi)=\min\{m,1-M\}\) & the factor in the trading bound \(\rho\,\kappa(\pi)\DT\)\\
\(\zeta(\bar P)\), \(e(x)\) & high-cost preparation probability; entry after order \(x\)\\
\(K_H\), \(K_L\) & gross return of a purchase in class \(H\) / a sale in class \(L\)\\
\(\Ent\) & entry: the probability that the challenger prepares\\
\(\Own_H\) & high-class ownership: the probability that a class-\(H\) challenger acquires the target\\
\(W_H(F)\) & probability that a prepared class-\(H\) challenger wins against \(F\)\\
\(r_0,r_1,r_2\) & weak, strong and very strong incumbent\\
\(r_C\) & the entry ceiling: \(\tau(r_C)=M(1/2)\)\\
\bottomrule
\end{tabularx}
\caption{Notation.}
\label{tab:02_model-notation}
\end{table}
